\section{Decision Accuracy, Transport, and Stopping: Proofs}\label{app:r19proofs}
This section gives the complete arguments for the new economic error chain.
The inherited verification, Bellman-bridge, and diffusion proofs are retained
in the preceding sections. The constants below pertain to the new stored
finite-law experiment unless a more general hypothesis is stated.

\subsection{Proof of Proposition~\ref{prop:r19decision}}
First suppose an optimal action $a^*$ exists. The approximate maximizing
property implies
\[
 r_\theta(s,a^*)+\widehat S_0(T_\theta(s,a^*))
 -r_\theta(s,\widehat a)-\widehat S_0(T_\theta(s,\widehat a))
 \leq\delta_\theta(s).
\]
Substituting $S_\theta=\widehat S_0-e+\Delta_\theta$ gives
\begin{align*}
 Q_\theta(s,a^*)-Q_\theta(s,\widehat a)
 \leq {}&\delta_\theta(s)
       +e(T_\theta(s,\widehat a))-e(T_\theta(s,a^*))\\
       &+\Delta_\theta(T_\theta(s,a^*))
                -\Delta_\theta(T_\theta(s,\widehat a)).
\end{align*}
Each difference is at most its action-set oscillation. For an unattained
finite supremum, take an $\epsilon$-maximizer and let $\epsilon$ decrease
to zero. This proves \eqref{eq:r19osc} without a differentiability assumption.

For finite actions let $v_j=e(T_\theta(s,a_j))-\bar e_\pi(s)$.
In the weighted inner product $\langle x,y\rangle_\pi=\sum_j\pi_jx_jy_j$,
the representer of $v_i-v_j$ has entries $1/\pi_i$ at $i$, $-1/\pi_j$
at $j$, and zero elsewhere. Its squared norm is $1/\pi_i+1/\pi_j$.
Cauchy--Schwarz gives
\[
 (v_i-v_j)^2\leq(1/\pi_i+1/\pi_j)\sum_k\pi_kv_k^2.
\]
Maximizing over the finite pair set proves
$\operatorname{osc}e\leq\sqrt{C_\pi R_\pi}$. The transport difference is
at most $2b_\theta$, proving \eqref{eq:r19riskloss}. For random fitted
objects the inequality is pathwise. With deterministic $C_\pi$, for example,
$\E\sqrt{C_\pi R_\pi}\leq\sqrt{C_\pi\E R_\pi}$ whenever the expectation
is finite. Repetition of the same random seed does not create a new draw
for this expectation. Finally, if every nonoptimal action loses at least
$g$ and the displayed bound is below $g$, a nonoptimal implemented action
would contradict the bound.\hfill$\square$

\subsection{Proof of Theorem~\ref{thm:r19transport}}
Couple both economies by the same innovation sequence and the same initial
postdecision state. If $\|X_k^\theta-X_k^0\|\leq d_k$, then the triangle
inequality and \eqref{eq:r19transportassumptions} give
\[
 \|X_{k+1}^\theta-X_{k+1}^0\|
 \leq\epsilon_{\theta,k}+L_kd_k=d_{k+1}.
\]
Induction starts at $d_0=0$. At each reward date,
$|g_{\theta,k}(X_k^\theta)-g_{0,k}(X_k^0)|
 \leq\eta_{\theta,k}+H_kd_k$.
Multiply by the nonnegative weights, sum including the terminal reward,
and take expectations. The bound is uniform because each assumed bound is
uniform on the invariant domain. No interchange of a supremum over future
controls with a fixed-policy expectation is used.

For recursive evaluations write
$D_k=\|V_{\theta,k}-V_{0,k}\|_\infty$. Add and subtract
$\mathcal G_{\theta,k}(V_{0,k+1})$ to obtain
$D_k\leq\nu_{\theta,k}+\beta_kD_{k+1}$.
Backward substitution from $D_K\leq b_K$ gives
\eqref{eq:r19recursive}. Both arguments require that every coupled state
and every recursively evaluated utility belong to the domains on which the
bounds were verified. The theorem makes no extension past these domains.
If a domain is localized, an exit contribution must be bounded separately
before using the result as a global guarantee.\hfill$\square$

\subsection{Proof of Proposition~\ref{prop:r19stopping}}
Strong concavity gives, for any feasible $b$,
\[
 Q(b)-Q(a)\leq Q'(a)(b-a)-\kappa(b-a)^2/2.
\]
In the interior, maximize this upper bound over all real $b-a$ to obtain
$|Q'(a)|^2/(2\kappa)$. At the lower endpoint restrict the maximization to $b-a\geq0$.
Its value is
\[
 \frac{(\max\{Q'(a),0\})^2}{2\kappa}.
\]
At the upper endpoint the analogous residual is $\max\{-Q'(a),0\}$. Enlarging a bounded
feasible half-interval to the half-line is conservative. This proves
\eqref{eq:r19residual}.

An outward derivative interval bounds the applicable residual, and a positive
lower curvature bound enlarges the resulting nonnegative fraction. Thus each
task certificate holds deterministically at every attempted stage. A first
successful average certificate implies the average economic target; previous
failures do not affect validity. If checks are stochastic, let $B_j$ denote
the event that attempted check $j$ falsely certifies its candidate. The
conditional hypothesis gives $\Pr(B_j)\leq\E\alpha_j$; an unattempted
check contributes zero. The union bound yields
$\Pr(\bigcup_jB_j)\leq\E\sum_j\alpha_j\leq\alpha$. This proof requires
conditional validity for the data actually used after past decisions.
A fixed-time interval repeatedly queried without such a property does not
meet the hypothesis.

The work identity follows by summing the disjoint measured operations through
the stop, including unsuccessful fits and tests. Under the stated exact
linear cost model, subtraction gives
$W_N(q)-W_R(q)=F_N-F_R-q(c_R-c_N)$, establishing the strict inequality
in \eqref{eq:r19break}. Integer query counts use that strict inequality;
a tie at an integer threshold is not a strict saving. Measured costs in the
experiment are used directly and are not assumed to satisfy a linear model.
\hfill$\square$

\subsection{Proof of Corollary~\ref{cor:r19amortization}}
Condition on the feature map and design, which are independent of the fitting
noises. Least squares gives
$\widehat\beta-\beta=G_n^{-1}\Phi^{\mathsf T}\xi$.
For a fixed queried argument $x$, the prediction error is a linear combination
of the independent noises. Its sub-Gaussian variance parameter is at most
\[
 \sigma^2\|\Phi G_n^{-1}\phi(x)\|^2
 =\sigma^2\phi(x)^{\mathsf T}G_n^{-1}\phi(x)
 \leq\sigma^2\ell/n.
\]
The two-sided sub-Gaussian bound and a union bound over the $K$ distinct
arguments therefore give
\[
 \Pr\left(\max_x|\widehat S(x)-S(x)|>
  \sigma\sqrt{\frac{2\ell}{n}\log\frac{4K}{\alpha}}\right)
 \leq\alpha/2.
\]
Uniformly on that event, the oscillation in every feasible finite action
menu is at most twice the displayed bound. Proposition~\ref{prop:r19decision}
with zero transport and \eqref{eq:r19learnbudget} gives loss at most
$\varepsilon$ at every query. Full column rank and the leverage bound are
separate hypotheses and must actually hold at the selected integer $n$.
They are not consequences of the scalar sample-size inequality.

For the direct method, each prediction average has sub-Gaussian variance
parameter at most $\sigma^2/t$. Dependence between predictions caused by
common innovations across different arguments does not affect the union
bound; independence is required only for the observations within an average,
or an equivalent sub-Gaussian bound for that average. The same calculation
with $\ell/n$ replaced by $1/t$ proves the second guarantee. A union bound
over the two methods gives joint coverage $1-\alpha$, without assuming
independence between them. Reusing a cached prediction at exactly the same
continuation identity and argument incurs no additional simulation count.
This is why $K$, rather than the number of repeated requests, enters the
work comparison.

Subtract the complete specified costs to obtain
\eqref{eq:r19explicitadvantage}. The fixed term $F(n)$ includes dictionary
construction, all labels, the least-squares solve and its validation; it is
not just a neural update clock. Common selection and verification work
cancels only under the explicit equal-cost hypothesis. Otherwise its
difference remains. These are sufficient budgets and actual costs for the
declared procedures, not statistical lower bounds on arbitrary algorithms.
If $|\E\widehat S(x)-S(x)|\leq B$ uniformly under a misspecified dictionary,
the prediction bound acquires $B$; if future continuation changes by at most
$b$, the decision bound acquires a further $2b$. Thus the same argument
uses margin $\varepsilon-\delta-2B-2b>0$, provided the stated bias bound is
verified. A pointwise approximation error of an arbitrary basis expansion
is not automatically a bound on the bias of its fitted least-squares
projection.\hfill$\square$

\subsection{Uniform curvature in the stored capital economy}
\label{app:r19curvature}
Let $M=(1-\mu)I+\mu\bm1\bm1^{\mathsf T}/d$. It is row stochastic and
has infinity operator norm one. For $f(y)=\lambda\tanh(My)$, the inequalities
$|\tanh'|\leq1$ and $|\tanh''|\leq2$ imply, along a scalar path with
$\|y'\|_\infty\leq P$ and $\|y''\|_\infty\leq D$,
\begin{equation}
 \|(f(y))''\|_\infty\leq\lambda(D+2P^2).
 \label{eq:r19prodcurv}
\end{equation}
The postdecision map has first derivative $-h\bm1$ and second derivative
zero, so initialize $P=h$, $D=0$. Adding the first innovation does not
change these derivatives. For each of the three future transitions set
\begin{align*}
 B_k&=\lambda(D+2P^2),&
 D&\leftarrow D+hB_k,&P&\leftarrow(1+h\lambda)P.
\end{align*}
The order matters: $B_k$ uses the preceding $P,D$. The second derivative of
the future production reward is at most $W_kB_k$, since every stored $W_k$
is positive. All future terms involving the fixed withdrawal are constant
with respect to the current action.

For a vector $v$ write $\operatorname{range}(v)=\max_i v_i-\min_i v_i$.
At the first shocked postdecision state a valid range bound for each stored
path is
\[
 R_0=\operatorname{range}(y)+2h\lambda+
                       \operatorname{range}(z_0).
\]
Every subsequent transition satisfies
$R_k=R_{k-1}+2h\lambda+\operatorname{range}(z_k)$ as an upper bound.
The scalar drift and common withdrawal cancel from a coordinate range.
For the terminal payoff $g(y)=-\gamma\operatorname{Var}_d(y)/2$,
\[
 (g(y))''=-\gamma\operatorname{Var}_d(y')
 -\gamma\overline{(y-\overline y\bm1)y''}
 \leq\gamma R_3D_3.
\]
The first term is nonpositive; dropping it is an upper bound, not a sign
reversal. Combining these inequalities gives a uniform upper bound on the
continuation curvature,
\[
 S_0(x(y,a))''\leq
 M(y):=\sum_{k=1}^{3}W_kB_k+\gamma\E_{16}[R_3D_3].
\]
Since the current reward has second derivative
$-A_0(w/a^2+\eta)$, a valid lower bound for true strong concavity on the
whole interval is
\begin{equation}
 \kappa(y,\theta)=A_0\{w/\bar a^2+\eta\}-M(y).
 \label{eq:r19kappa}
\end{equation}
The verifier evaluates this expression with outward arithmetic and requires
a strictly positive lower endpoint for every task. A nonpositive result
would be a verification failure, not a reason to delete that task. The
published records retain every returned lower endpoint.

The true derivative is enclosed separately rather than inferred from the
surrogate. For each stored path initialize $p_0=-h\bm1$. At a future state,
\[
 f'(y)=\lambda\{1-\tanh^2(My)\}Mp,\qquad
 p\leftarrow p+h f'(y).
\]
Accumulate $W_k\overline{f'(y)}$ and the terminal derivative
$-\gamma\overline{(y-\overline y\bm1)p}$. Average over all sixteen atoms
and add the current derivative
$A_0(w/a-\eta a-\tau)-W_0$. Every operation is evaluated on enclosing
intervals. Dependency between occurrences of the same variable may widen
an interval, but cannot invalidate the enclosure.

\subsection{Arithmetic enclosure and its implementation}
The stored binary64 coefficients and innovation increments define the exact
finite economy. Floating-point evaluation of a Gaussian generator is not
being asserted to be exact sampling from a continuous Gaussian economy.
The generator is used to construct and freeze a finite array; the subsequent
scientific statement concerns that array with equal masses.

The arithmetic module encloses each binary operation by the adjacent
representable numbers toward minus and plus infinity. It checks finite
values, IEEE round-to-nearest behavior and gradual underflow. Products and
quotients take extrema over endpoint combinations, with divisors bounded
away from zero. Squaring includes zero when the argument interval contains
zero. Balanced reductions avoid interpreting a non-enclosing ordinary mean
as an interval sum. Induction on an expression's operation graph proves
inclusion for arithmetic expressions from enclosed inputs.

The nonlinear enclosure does not assume that a library's hyperbolic tangent
is correctly rounded. For $0\leq t\leq1/2$, the exponential is enclosed by
its degree-18 Taylor polynomial and the nonnegative tail bound
\[
 0\leq e^t-\sum_{j=0}^{18}t^j/j!
 \leq \frac{t^{19}/19!}{1-t/20}.
\]
All terms are enclosed using the same arithmetic. Exact binary scaling (checked on every replayed argument) reduces the
argument to this interval; repeated squaring and reciprocation recover
positive and negative arguments. The identity
$\tanh x=(e^{2x}-1)/(e^{2x}+1)$ and monotonicity provide endpoint enclosures.
For the auxiliary logarithm check, scaling to $[1,2)$ gives
$z=(x-1)/(x+1)\in[0,1/3)$ and
\[
 \log x=2\sum_{j=0}^{31}\frac{z^{2j+1}}{2j+1}+R,
 \qquad 0\leq R\leq\frac{2z^{65}}{65(1-z^2)}.
\]
The stored enclosing endpoints for $\log2$ is independently checked with
exact rational series bounds. Synthetic high-precision derivative tests
check the implemented recursion on separate small-dimensional cases.
These tests complement the analytical enclosure argument; they are not
substitutes for it or additional economic observations.

\subsection{Charge monotonicity and the economic tolerance}
For an interior maximizing action, differentiate
$Q_\theta'(y,a^*(\tau))=0$. Because the charge enters the derivative as
$-A_0\tau$, the implicit-function identity is
$Q''\,\partial a^*/\partial\tau-A_0=0$, proving
\eqref{eq:r19charge}. For the constrained case take $\tau_2>\tau_1$ and
respective optimal actions $a_2,a_1$. Adding their optimality inequalities
leaves
$A_0(\tau_2-\tau_1)(a_2-a_1)\leq0$, hence $a_2\leq a_1$.
Uniqueness follows from strict concavity.

An external grant in the utility account satisfies
$A_0w\log(1+g)=U$ exactly when $g=\exp(U/(A_0w))-1$. This is a
monotone normalization of a certified payoff loss. The published grant
table uses ordinary binary64 evaluation and is explicitly descriptive;
the primary rigorous statement is the outward payoff-loss bound. No claim
is made that increasing actual endogenous consumption by this amount leaves
capital, future rewards, or the reference policy unchanged.
